\cotefile{123}
\scanpage{37}
% The curved lead runs from between « remarquables » and « et » to the two lighter lines.
Liste d'espaces homogènes remarquables\inserted{liés aux ss-groupes simples \unsure{principaux} de rang $1$,} et leurs relations
% archive sheet number: 36

\scanpage{38}
\uline{Schémas} :

\diagram{\scriptsize
\(\left\{\begin{array}{ll}
\text{des sousgroupes simples \inserted{(de rang $1$)} principaux} & \mathrm{SSP}\\
\text{--- , \uline{avec épinglage}} & \mathrm{SSPEP}\subset\underline{\mathrm{Hom}}_{\mathrm{gr}}^{*}(\mathrm{S}\unsure{\ell|L|l}(2),G)\\
\text{--- , avec couple de Kill} & \mathrm{SSP}\deleted{\unsure{\mathrm{EP}}}\mathrm{Kill}\\
\text{--- , avec tore max} & \mathrm{SSPTor}\\
\text{--- , avec tore maximal ``pointé''} \ \unsure{\simeq|\sim} G_{\mathrm m} & \mathrm{SSPTor}_{0}\subset\underline{\mathrm{Hom}}_{\mathrm{gr}}(G_{\mathrm m},G)\\
\text{--- , avec Borel} & \mathrm{SSPBor}\\
\text{--- , avec Borel \unsure{éping.}} & \mathrm{SSPBor}.\mathrm{ep}
\end{array}\right.\)}
\margin{\deleted{\unsure{et désignons leurs intersections}}}

\diagram{\small
\(\left\{\begin{array}{ll}
\text{des tores (\deleted{\unsure{de}}dim $1$) principaux} & \mathrm{STP}\\
\begin{array}{l}\text{des hom }G_{\mathrm m}\unsure{-|\to}G\text{ principaux}\\
\quad\text{(i.e. sous-tores ``épinglés'')}\end{array} & \mathrm{ST}_{0}\mathrm P
\end{array}\right.\)}

\diagram{\scriptsize
\(\left\{\begin{array}{ll}
\text{des sous-groupes additifs principaux} & \mathrm{SAP}\\
\begin{array}{l}\text{des sous-groupes additifs principaux}\\ \qquad\text{pointés}\end{array}
 & \mathrm{SA}_{0}\mathrm P\subset\underline{\mathrm{Hom}}_{\mathrm{gr}}(G_{\mathrm a},G)\\
\text{des \deleted{éléments} \unsure{unipotents|unip. tout} principaux} & \mathrm{ElUP}\subset G
\end{array}\right.\)}
\margin{$\big\}$ sont isomorphes can. en car. 0}

\deleted{des éléments} \uline{NB} \deleted{\illegible{1}} \unsure{peut} donner des variantes en termes de $\unsure{\mathfrak g|V}$ des trois derniers cas\inserted{(et sans doute d'autres)}, qui \deleted{\unsure{leur}} sont respectivement can. isom. en car. 0 tout au moins.

\noindent\rule{\linewidth}{0.3pt}

% Three oblique cancellation strokes pass through the catalogue below the rule.
% Local cancellations beneath that overall cancellation are marked inside it.
\deleted{des éléments \unsure{périodiques|réguliers} principaux \deleted{\illegible{1}}}\hfill
$\deleted{\mathrm{ECox}}\quad\deleted{\unsure{\mathrm{ECox}'}}$

\deleted{(\uline{NB} \unsure{au cas quelle|au cas quelle,}\inserted{en tous cas sur \unsure{Alg.} clos,} il vaut mieux remplacer par les hom principaux $\mu_h \unsure{-|\to}G$ ?)}\hfill
$\deleted{\mathrm{Kost}}\ \deleted{\unsure{\mathrm{Kost}_{h}}}$

\margin{\deleted{$\big\}$ deviennent isomorphes une fois choisie une racine primitive $h$.\textsuperscript{ème} de 1 (dans \unsure{son} $\mu_h^{*}$ !)}}

\diagram{\scriptsize
\(\left\{\begin{array}{ll}
\deleted{\text{des droites cycliques}\ \subset\mathfrak g} & \deleted{\mathrm{DC}}\\
\deleted{\text{des \deleted{\unsure{droites}} éléments cycliques}\ \in\mathfrak g} &
\deleted{\mathrm{DC}_{0}\simeq\mathrm{EC}}\\
\deleted{\text{des éléments \deleted{principaux}\inserted{cycliques} principaux}} &
\deleted{\mathrm{ECP}}\\
\end{array}\right.\)}
\margin{\deleted{\uline{NB} ce n'est pas un esp. homogène $G$ ; on a $\mathrm{DC}_{0}/G\simeq E^{1}\deleted{\unsure{G_{\mathrm m}}}$ (si on a une droite cyclique\inserted{$D$}, l'action de son normalisateur\inserted{dans $G$} \unsure{se fait} via un groupe cyclique fini\inserted{$\unsure{\Gamma\subset k^{\times}}$}, \unsure{à valeurs|et valeurs} \unsure{multipli|multipliée}\illegible{1} de là, \unsure{pr. le}\illegible{1} \deleted{\unsure{supprimé}} c'est la \unsure{quantité considérée} qui est $D/\Gamma$)}}

\diagram{\scriptsize
\(\left.\begin{array}{ll}
\deleted{\text{des couples de }\deleted{\unsure{\text{Coxeter}}}\inserted{\unsure{\text{Co}\deleted{\text{standard}}+\text{Co}^{\prime}\text{élect.}}}\ (T,\sigma\in N(T))}
 & \deleted{\unsure{\mathrm{CoxCo}k_{0}}}\ \deleted{\unsure{\mathrm{Cox}_{0}|\mathrm{Clox}}}\\
\deleted{\text{des couples de tores maximaux en opposition}} & \deleted{\mathrm{CTap}}\\
\left\{\begin{array}{l}\deleted{\text{des couples de Coxeter }(T,c\in W(T))}\\[3pt]
\deleted{\text{des couples de Kostant }(\kappa:\mu_h\unsure{-|\to}G,D)}\end{array}\right.
& \begin{array}{l}\deleted{\mathrm{CCox}}\\[3pt]\deleted{\mathrm{CKost}}\end{array}
\end{array}\right\}\ \deleted{\text{sont isom sur }\mu_{h}^{*}}\)}
% A curved lead points up from the Coxeter row to the first « des couples » row.
% archive sheet number: 37

\scanpage{39}
% Very pale inverted type shows through the sheet from its other face.
\diagram{\scriptsize
\begin{tikzcd}[ampersand replacement=\&,column sep=0.55em,row sep=0.55em]
 \& \& \& \underline{\mathrm{Isom}}(G_0,G)
     \arrow[dr,dashed] \arrow[dd,phantom,"\simeq" description,sloped] \& \& \\
 \& \& \& \& \underline{\mathrm{Hom}}_{\mathrm{gr}}(\mathrm{S}\unsure{\ell|L|l}(2),G) \& \\
 \& \& \& \mathrm{SSPEP}\arrow[ur,phantom,"\unsure{\subset|\hookrightarrow}" description,sloped]
    \arrow[dl,dash]\arrow[dd,dash]\arrow[dr,dash] \& \& \\
 \& \mathrm{Kill}\arrow[d,phantom,"\simeq" description,sloped] \&
 \mathrm{SSPTor}_{0}\arrow[dl,dash]\arrow[dd,dash]
 \arrow[dr,phantom,"\simeq" description,sloped] \& \&
 \mathrm{SSPBor}\unsure{.\mathrm{ep}|\mathrm{ep}}\arrow[dd,dash]\arrow[dr,dash] \& \\
 \& \mathrm{ST}_{0}\mathrm P\arrow[dd,dash]\arrow[dl,dash,very thick,"\unsure{?}"'] \& \&
 \mathrm{SSP}\deleted{\unsure{\mathrm E}}\mathrm{Kill}\arrow[dl,dash]\arrow[dr,dash] \& \&
 \mathrm{SA}_{0}\mathrm P\unsure{.}\simeq\mathrm{ElUP}\arrow[dd,dash]\\
 \deleted{\unsure{\mathrm{ECox}_{h}}}\mathrm{Kost}\arrow[dd,dash,very thick] \& \&
 \mathrm{SSPTor}\arrow[dl,dash]\arrow[dr,dash] \& \&
 \mathrm{SSPBor}\arrow[dl,dash]\arrow[dr,dash] \& \\
 \& \mathrm{ST}\deleted{\unsure{0}}\mathrm P\arrow[dl,dash] \& \& \mathrm{SSP} \& \&
 \mathrm{SAP}\deleted{\unsure{\mathrm{ElUP}}}\\
 \mathrm{Tor} \& \& \& \& \&
\end{tikzcd}}
\margin{diagrammes d'\unsure{espaces} homogènes sous $G$.}

\diagram{\small
\begin{tikzcd}[ampersand replacement=\&,column sep=0.8em,row sep=1.2em]
 \& \& \& e\arrow[dl,dash,"\underline{1}"']\arrow[dd,dash,"\underline{1}"]\arrow[dr,dash,"\underline{1}"]
 \& \& \\
 \& \& t\arrow[dl,dash,"r-1"']\arrow[dd,dash,very thick,"{\mathbb Z/2\mathbb Z\deleted{\unsure{0}}}"]
 \arrow[dr,phantom,"=" description,sloped] \& \&
 b_u\arrow[dd,dash,"\underline{1}"]\arrow[dr,dash,"\underline{r-1}"] \& \\
 \& T\arrow[dl,dash,"0"']\arrow[dd,dash,"{\mathbb Z/2\mathbb Z}"] \& \&
 \mathbf{t}\arrow[dl,dash]\arrow[dr,dash,"\underline{1}"] \& \&
 Z(b_u)\arrow[dd,dash,"{\underline{1}\,(G_{\mathrm m})}"]\\
 ?\ \unsure{N|\mathbb N}_1\arrow[dd,dash] \& \&
 \unsure{n|\mathfrak n}\arrow[dl,dash,"r-1"']\arrow[dr,dash,"\underline{2}"] \& \&
 b\arrow[dl,dash,"\underline{1}"']\arrow[dr,dash,"\underline{r-1}"] \& \\
 \& N_0\arrow[dl,dash] \& \& H \& \& N(b_u)\\
 N \& \& \& \& \&
\end{tikzcd}}
% archive sheet number: 38

\scanpage{41}
\uline{Résultats de Kostant} sur classes de conjugaison remarquables dans $G$ $1/2$ simple adjoint (sur $\mathbb C$)

a) Classes de conjugaison de ss-groupes simples de rang $1$ $\deleted{\unsure{\simeq}}\unsure{\simeq}$ classes de conjugaison d'éléments unipotents de $G$ [ou de $\unsure{\mathfrak g|V}$\inserted{ou de ss-groupes additifs de $G$},] $\simeq$ certaines classes de conjugaison \deleted{d'éléments semi-simples de $\mathfrak g$ [\uline{NB} \unsure{en de hom}\illegible{1} \unsure{princ.} ss} d'homomorphismes $\mathbb G_{\mathrm m}\to G$ [ou d'éléments de $\mathfrak g$ \ldots].

Les hom admissibles $G_{\mathrm m}\to T\simeq\deleted{\unsure{T^{\Sigma}}}G_{\mathrm m}^{\Sigma}$ ($\Sigma$ ensemble des racines simples) dans le ``domaine fondamental'' $W$ de $\mathrm{Hom}(G_{\mathrm m},T)\simeq\mathbb Z^{\Sigma}$, savoir $\mathbb N^{\Sigma}$, correspondent à \inserted{des} éléments dont les composantes sont $0,1,2$ (pas tous)

% margin note written sideways
\margin{l'élément $w_0\in W$ qui transforme la chambre $C$ en l'opposée}

b) Le cas où toutes les composantes sont $1$ est admissible ; il correspond aux \uline{éléments unipotents principaux}\inserted{(ou \uline{réguliers})}, et dans $\mathfrak g$ aux \uline{éléments semi-simples principaux} ; \deleted{\unsure{aux}} \uline{hom} ``\uline{principaux}'' $\mathbb G_{\mathrm m}\to G$. On parle des sous-groupes simples de rang $1$ \uline{principaux}.

\deleted{\unsure{b)}}

Centralisateur d'un élément unipotent principal : sous-groupe \uline{connexe}\inserted{(\unsure{Springer|de Springer})} \inserted{\uline{unipotent}} \uline{commutatif}\inserted{\uline{de dim $\unsure{r}$}}

$\times$\quad Normalisateur d'un ss-groupe additif principal : groupe connexe \deleted{\unsure{de dim}}\inserted{\unsure{son quotient} par \unsure{précéd.} $\simeq G_{\mathrm m}$ (\unsure{canont})}

Centralisateur d'un hom principal $\mathbb G_{\mathrm m}\to G$ : un tore maximal. Donc ces hom sont en correspondance biunivoque avec les couples de Killing de $G$.

\margin{Donc sur base quelc $\exists$ ss-\deleted{\unsure{groupe simple}}\inserted{de} \deleted{\unsure{de}} rg $1$ princip. \uline{déployé} ssi $G$ quasi-déployé}

Normalisateur d'un $1$-tore principal : groupe $\deleted{\unsure{N}}N_0$ dans lequel $T$ est d'indice $2$, correspondant à \unsure{$b|w$}

\deleted{\unsure{Centralisateur}}\inserted{Normalisateur} d'un sous-groupe \deleted{principal} de rang $1$ principal : lui-même.

\deleted{Normalisateur}\inserted{Stabilisateur} d'un ss-gpe principal de rang $1$\inserted{$H$}, avec un tore maximal $T_0$ donné dans celui-ci (resp. un \deleted{couple de}\inserted{Borel, resp. un} Killing) : Normalisateur\inserted{$N_0$} du tore\inserted{$T_0$} dans $H$\unsure{.} (resp. $B_0$, resp. $T_0$)

\margin{\deleted{\unsure{Donc sur un corps} \illegible{1} couple \unsure{de rang} $1$}\inserted{\uline{NB} $\mathbb Z/2\mathbb Z$ opère sur \unsure{KillCl}, échangeant les couples opposés.} \deleted{\unsure{là-dessus, on a}\illegible{2}} \unsure{quotient est le schéma des} \deleted{\unsure{ss}} tores de dim $1$ \uline{principaux}}

c) Soit $G_{\mathrm m}\xrightarrow{u}G$ un hom principal, et soit \deleted{\illegible{1}} $h$\inserted{$=o(\psi)+1$} \deleted{\unsure{le}} le nb de Coxeter ($=$ somme des coeff. de la plus haute racine $\psi$, \uline{plus} $1$).

Soit $\zeta$ \deleted{\unsure{un él.|une}} racine primitive $h$.\textsuperscript{ème} de $1$, et considérons $u(\zeta)$. On prouve que les $u(\zeta)$ ($\zeta$ variable, et évidemment aussi pour $u$ variable) sont conjugués. Les éléments obtenus sont appelés \inserted{\uline{éléments}} \uline{principaux} \inserted{\uline{réguliers}} de $G$. Si $T$ est un tore maximal, et $\sigma\in W=N(T)/T$ \deleted{\unsure{élémt}} un \deleted{\unsure{élémt}} élément de Coxeter, on prouve que les représentants de $\sigma$ sont tous principaux (donc réguliers).

\margin{\unsure{en car qui|encore que|ordre qui} divise $h$, il vaut mieux considérer la classe d'hom. $\mu_h\to G$}

Centralisateur d'un élément régulier principal $\gamma$ : ss-gpe de $N(T)$ \deleted{\unsure{formé}} d'image réc.\ du ss-gpe de $W=W(T)$ qui stabilise $\gamma$.

\margin{\textcircled{?}} C'est probablement $T$, i.e. \deleted{\unsure{on a ainsi}} la donnée de $\gamma$ revient à celle de $\unsure{\varphi|u|\mu}:G_{\mathrm m}\to G$

\margin{$G$ \uline{simple}}
d) Considérons l'opération de $W$ sur $\deleted{\unsure{\mathfrak g}}$ $\underline t=\mathrm{Lie}(T)$. Il existe des \deleted{\unsure{dr. propres}}\inserted{droites} propres\inserted{$D$} \deleted{\unsure{pour|dans}} $\sigma$ (op. de Coxeter) dont la valeur propre est une racine primitive\inserted{$h$.\textsuperscript{ème}} de $1$, \inserted{ce sont exactement celles qui sont ``régulières'',} et ces droites sont conjuguées sous l'action de $W$ (donc : \unsure{fusion|faisceau|forment} conjuguées
% archive sheet number: 40

\scanpage{42}
dans $\mathfrak g$ sous $G$) ; \deleted{\unsure{de même|donc}, les droites} ``plus généralement, si $g$ est un élément \deleted{régulier principal} de $G$, des \deleted{\unsure{toutes}}\inserted{ses} droites propres\inserted{\uline{régulières}} sous $g$ de $\mathfrak g$ de val. propre une racine prim. $h$.\textsuperscript{ème} existent \deleted{ssi} $g$ est régulier principal, et les droites en question sont conjuguées. On \deleted{\unsure{les}} appelle de telles droites de $\mathfrak g$ cycliques, les pts correspondants ($\ne 0$) de $\mathfrak g$ \uline{cycliques}. \deleted{\unsure{F}} On en construit, avec un couple de Killing donné, en posant $\sum_{1}^{\unsure{r|n}}\lambda_i e_{\alpha_i}+\mu e_{-\psi}$, les $\lambda_i$ et $\mu$ tous $\ne 0$, ($\psi$ la plus haute racine). \unsure{$\Vert$}
% margin note written sideways
\margin{\deleted{\unsure{Correspondant à} tout \unsure{Tore|ss-tore} de $\unsure{G|C}$, \deleted{\unsure{qui est}}\inserted{il} $\exists$\unsure{ra} droite \deleted{\unsure{cyclique}}\inserted{une} stable par $T$ \unsure{d'où un}}}

Le centralisateur d'un élément cyclique\inserted{$x$} de $\mathfrak g$ est un tore maximal $T$, \deleted{le centralisateur est un groupe $\deleted{\unsure{T^{\prime}_{0}}}T'$ entre $T$ et $N(T)$ (ou plutôt il s'agit de la droite $D$ engendrée par $x$) ;} le normalisateur de $D$ dans $G$ est un groupe $T^{\unsure{\sim|\prime}}$, $T\subset T^{\sim}\deleted{\unsure{G'}}\subset N(T)$, $T^{\sim}$ contient l'élément de Coxeter, donc \uline{n'est pas réduit à} $T$. Il faut \deleted{\unsure{donc}} déterminer les ss-groupes correspondants $\deleted{\unsure{W'}} W^{\sim}$ de $W$ ($\deleted{\unsure{W'}}W^{\sim}\subset C^{*}$ est un groupe cyclique\inserted{d'ordre $h'$ multiple de $h$,}\deleted{\unsure{on peut}} \inserted{\unsure{donc en|donc égal à} $h$, grâce à la propriété maximale de l'élément de Coxeter, de \unsure{$T$}.} \deleted{\unsure{on peut voir que}} il est engendré par l'élément de Coxeter. On voit que si $D$ est une droite critique, alors $\unsure{\{}D/\mu_{h'}\simeq E^{1}$ canoniquement [\unsure{se réduire|se rédui|se vérifie}\illegible{1} au cas quasi-déployé, et dans ce cas on a un élément privilégié cyclique, savoir $\sum_i e_{\alpha_i}+e_{-\psi}$. Donc on peut regarder les él\textsuperscript{ts} cycliques qui donnent $1$ dans $\deleted{\unsure{K\simeq|E}}\mathrm{EC}/G\simeq E^{1}$, on peut les appeler les éléments \uline{cycliques principaux} (dépendent à priori du choix d'un signe, \deleted{si $x$ est principal, alors $\unsure{\zeta}x$ est \unsure{antiprincipal}, où $\unsure{\zeta}$ est une racine primitive $2h$.\textsuperscript{ème} de $1$}). Les él\textsuperscript{mts} forment un torseur sous $G$, de stabilisateur un tore maximal $T'$. \deleted{Dans le groupe déployé, \unsure{cet}}\inserted{Pour \deleted{\unsure{que}} un} élément de la forme $\sum e_{\alpha_i}-e_{-\psi}$, on trouve
\margin{\uline{NB} \deleted{\unsure{Le}} le centralisateur d'un sous-tore contenant un élément régulier est un tore maximal}
\margin{\textcircled{?}}
\margin{Mais peut-on choisir un gén. can. de $e_{-\psi}$ ? ici le gpe. des autom \deleted{\unsure{de ce fibré}} d'un gpe épinglé (respectant l'épinglage --) induit-il l'identité sur $P_{-\psi}$ ? \deleted{\illegible{1}} ou mieux sur $P_{\pm\psi}$ ? Sinon, peut-il \ldots}
\margin{préciser comment ?}

e) Considérons les couples $(T,\unsure{g|\gamma})$ d'un tore maximal de $\unsure{G,|G/}$ et d'un élément $g\in N(T)$ \deleted{dont} l'image dans $W$ est un élément de Coxeter. Deux tels couples sont conjugués [\deleted{\unsure{donc}} si on a $(T,g')$ un autre tel couple, $\exists\,n\in N(T)$ tel que $ngn^{-1}=g'$]. Le sous-groupe de $G$ qui stabilise un couple $(T,g)$ est \deleted{discret} un sous-groupe \uline{discret} de $N(T)$, contenant $g$, c'est $N^{\unsure{\sigma|\gamma}}$, extension de $W^{\sigma}$ (car $\mathrm H^1(\mathbb Z/h\mathbb Z,T)=0$) par $T^{\sigma}\deleted{\unsure{est}}\simeq P/(1-\sigma)P\,(1)$ \unsure{--- twist Tate}.
% archive sheet number: 41

\scanpage{43}
Résumé de quelques résultats de Kostant
% archive sheet number: 42

\scanpage{44}
$G$, groupe semi-simple adjoint sur base $S$, de type constant.

(i) À tout couple de Killing $(T,B)$ dans $G$, on associe un hom. (en fait un mono)
\[
G_{\mathrm m}\overset{i}{\lhook\joinrel\longrightarrow}G
\]
par descente à partir du cas où le système de racines sur $T$ est déployé, donc où $T\xrightarrow{\sim}G_{\mathrm m}^{\Sigma}$ ($\Sigma$ ens des racines simples), \deleted{\unsure{on prend|quand}} d'où $\mathrm{Hom}(G_{\mathrm m},G_{\mathrm m})\simeq\mathbb Z^{\Sigma}$ : on prend l'élément dont toutes les coordonnées sont $1$.

Le centralisateur d'un tel hom (i.e. d'un tel sous-tore) \deleted{est} $T$, donc le $\unsure{=|\sim}$ que pour le couple de Killing.

On appelle \uline{hom principal} $i$ de $G_{\mathrm m}$ dans $G$ un hom tel qu'il existe un couple de Killing auquel il est associé. Ce couple est alors unique [$i$ est régulier], et par suite la condition envisagée sur $i$ est loc\textsuperscript{le} pour fpqc. Il y a corr. biunivoque entre hom. principaux $G_{\mathrm m}\unsure{-|\to}G$ et couples de Killing.

(i bis) \deleted{Le sous-tore normalisateur de $i(G_{\mathrm m})$, $i$ principal, est} Si $i$ est principal, $i^{-1}=[\lambda\mapsto i(\lambda^{-1})]$ aussi, il est associé au couple de Killing opposé. Le normalisateur de $i(G_{\mathrm m})$ est le ss-groupe\inserted{$N_{\unsure{1|0}}(T)$} de $N(T)$ image inverse du sous-groupe $\simeq(\mathbb Z/2\mathbb Z)_S\subset W$ engendré par l'élément $w_{(T,B)}$ de $W$ qui envoie $B$ sur le Borel opposé.
\margin{$r\geq 1$ !}

Sous-tore de rang $1$ principal : qui loc. fpqc est \deleted{\unsure{de}}l'image d'un hom principal $\Longleftrightarrow$ qui sur un rev. étale \deleted{\unsure{galoisien}} principal de \deleted{groupe} $\mathbb Z/2\mathbb Z$ est comme \unsure{ça}.
% archive sheet number: 43

\scanpage{45}
Il y a correspondance biunivoque entre \uline{ces} sous-tores et les sections de $\mathrm{Kill}/(\mathbb Z/2\mathbb Z)$ [NB ``\uline{le}'' $W$ opère sur $\mathrm{Kill}$, et dans $W$ il y a ``\uline{le}'' élément $\underline w$ d'ordre $2$ qui transforme \uline{la} \deleted{B} chambre de Weyl\inserted{\unsure{standart|standard}} en son opposée, c'est par ce $\underline w$ qu'on divise \ldots] : couples de un tore maximal et d'un couple de Borel opposés le contenant (mais ces Borels n'étant pas néc. définis séparément sur $S$).

\margin{\unsure{si|G} simple ?}
(ii) Soit $i:G_{\mathrm m}\unsure{-|\to}G$ un hom principal. Il induit un hom
\[
\unsure{\kappa|K}_{\deleted{\unsure{h}}}:\mu_h\longrightarrow G.\qquad h\ \text{le nb de Coxeter}
\]
Un hom $\kappa_{\deleted{\unsure{h}}}$ est dit \deleted{\unsure{de}}\inserted{\uline{hom}} \uline{de Kostant} si loc fpqc (ou ét.f., cela vient au $=$, comme on voit aussitôt) il est obtenu à partir d'un hom principal.

\deleted{Le} Alors $\kappa_{\deleted{\unsure{h}}}$ est régulier, i.e. le centralisateur connexe de $\kappa_h$ est un tore maximal. Faisant opérer $\mu_h$ sur $\mathfrak g\mathrel{\underset{\unsure{*}}{=}}\underline{\mathrm{Lie}}(G)$, la graduation de type $\mathbb Z/h\mathbb Z$ correspondante de $\mathfrak g$ donne comme composante de degré $0$ une \deleted{\unsure{tore maximal} $T=$ \unsure{Zer}}\inserted{sous-Algèbre de} Cartan $\mathfrak t=\deleted{\unsure{\mathfrak g^0}}\mathrm{Lie}(T)$, $T=\underline{\mathrm{Cent}}^{0}(i_h)$, et la composante de degré $j$, $1\leq j\leq h-1$, \deleted{est égal au}\inserted{a comme} rang le nombre de racines positives dont l'ordre est $\equiv j\ (h)$, i.e. \unsure{don|dont} l'ordre est $j$ ou $j-h$ ; en fait, si $\kappa$ est associé à un couple de Killing, \deleted{\unsure{sur}}, \deleted{déployé} avec syst. de racines associé à $T$ \uline{déployé}, alors
% archive sheet number: 44

\scanpage{46}
la composante ${\mathfrak g}^{j}_{\deleted{\unsure{\kappa}}}$ ($1\leq j\leq h$) de $\mathfrak g$ suivant le caractère $\chi_j$ est
\[
\sum_{\substack{\alpha\in R\\o(\alpha)=j}}\mathfrak g^{\deleted{\unsure{j}}}_{\alpha}
+\sum_{\substack{\alpha\in R\\o(\alpha)=j-h}}\unsure{\widetilde{\mathfrak g}|\mathfrak g}_{\alpha}.
\]
En particulier $\mathfrak g^{1}$ est de rang $r+1$ ($r=$ rang).

Le centralisateur\inserted{$G^{\unsure{\kappa|K}}=N(T)^{\kappa}$} de $\kappa_{\unsure{h}}$ est un sous-groupe lisse de $N(T)$, image inverse du sous-groupe $\deleted{\unsure{W'|W}}$ de $W_{\kappa}$, \deleted{centralisateur} stabilisateur de $\kappa$ dans $W$. Kostant dit que le rang de $W_{\kappa}$ est égal à celui du centre de $\widetilde G$, et on peut espérer qu'il soit isomorphe à $\overline{P/R}$, $P=\widehat T=$ réseau des poids, $R=$ réseau des racines, \deleted{\unsure{ou}} ou à son dual \deleted{\unsure{de T}} à valeurs \unsure{ds|sur|dans} $\mathbb Q/\mathbb Z$.

La donnée d'un
\[
\kappa:\mu_h\unsure{-|\to}T
\]
de Kostant dans un tore maximal définit sur le système\inserted{de racines $R(T)$ de ce} dernier un système de racines \unsure{Q-simple|Qu-simple|Q-simples}\inserted{$\Sigma'$}, ou \unsure{Q-base|Qu-base} : l'ens des racines $\alpha$ telles que $\mathfrak g_{\alpha}\subset\mathfrak g^1$ ($\mathfrak g^1$ désigne la composante de $\mathfrak g$ suivant le car. $\chi_1$ de $\mu_h$). \inserted{En fait,} $\kappa$ est connu quand on connaît $\deleted{\unsure{Q}}\Sigma'\subset R$. \deleted{\unsure{Pb} des signes si $r>1$ d'où $h>2$ (donc la symétrie de $\mu_h$ n'est pas l'identité)} Il revient au $=$, dans le cas déployé, de dire que le \deleted{centralisateur} stabilisateur\inserted{$W_{\kappa}$} de $\kappa$ dans $W$ est \uline{égal} (à priori il est $\subset$) le stabilisateur de $\Sigma'$ dans $W$, i.e. l'ens des $w\in W$ tels que $w(\Sigma)\subset\Sigma'$, i.e. tels que $w$ transforme \uline{toutes} les racines de $\Sigma$ en
% archive sheet number: 45

\scanpage{47}
racines de $\Sigma$, sauf une qui est transformée en $-\psi$ ($\psi=$ plus grande racine).

Déterminer aussi le normalisateur de $\unsure{\kappa|K}$, et l'action de $\mathrm{Norm}(\kappa)/\mathrm{Zent}(\kappa)$ sur $\mu^{\unsure{\circ}}_h$ ; trouve-t-on tout $\mathrm{Aut}(\mu_h)=(\mathbb Z/h\mathbb Z)^{*}$ ?

\margin{$G$ simple ?}
(iii) Soit $\deleted{\unsure{\kappa}}$
\[
\kappa:\mu_h\unsure{-|\to}T
\]
un hom principal. Considérons comme dessus ${\mathfrak g}^{1}\inserted{=\sum_{\alpha\in Q}\mathfrak g^{\alpha}\ \text{(qui détermine $\kappa$)}}$. Kostant prouve (en car. $0$)

a) Toute section $\sum x^{\alpha}\in\Gamma\mathfrak g^1$ $(\unsure{x^{2}|x^{\alpha}}\in\Gamma\mathfrak g^{\alpha}]$ telle que les $x^{\alpha}$ ne s'annulent pas est \uline{$1/2$ simple régulière}.
% slanted
\margin{pourrait marcher au mieux si $h$ est inv. sur $S$}

b) Si un des $x^{\alpha}$ est nulle, elle est \uline{nilpotente}.
% slanted
\margin{marche sans doute $\overline{\mathrm{car}}$. Elle est nilpotente régulière ssi tous les autres ne s'annulent pas ?}

c) Soient deux droites $D,D'\subset\check V(\mathfrak g^{1})^{\mathrm{ss}}$, qui soient du cas a). Alors $\exists$\inserted{loc. fpqc} $t\in\Gamma T$ telle que $D'=\mathrm{ad}(t).D$

[Sans doute le stabilisateur dans $T$ d'un $D$ comme dessus est $\kappa(\mu_h)$, donc $t$ est une section canonique de $T/\kappa(\mu_h)$. On peut dire aussi que \deleted{\unsure{tout forme les isom.}}\inserted{dans} $\underline{\mathrm{Isom}}_{\underline{0}}(D\overset{\sim}{\to}D')$ \deleted{\unsure{qui envoient} \illegible{1}, \unsure{classe de}\illegible{1}}\inserted{on peut définir, un sous-}\deleted{\unsure{fibré}} $\mu_h$-torseur canonique [savoir \deleted{\unsure{l'ens des}} le fibré des $t\in T$ qui envoient $D$ dans $D'$], et pour un isom donné de $D$ avec $D'$ dans cette classe, il $\exists 1$ (unique) $t\in T$ qui l'induit [Cette information résulte du fait que un élément de $T$ qui est l'identité sur $D$ est la section $1$].
% archive sheet number: 46

\scanpage{48}
(iv) \uline{Droite} cyclique $D\subset\check V(\mathfrak g)$ : qui loc. fppf est propre de caractère $\chi_1$ sous $\mu_h\overset{\unsure{\kappa|K}}{\unsure{-|\longrightarrow}}G$, $\kappa$ \unsure{étant} de Kostant. Les droites cycliques sont donc conjuguées. Si $D$ est une droite cyclique, alors (sous réserve que a) \unsure{marche|unique}) $\mathrm{Zent}^{o}(D)$ est un tore maximal, $T'$, et $\mathrm{Norm}(D)$ est un sous-gpe de $N'=N(T')$ contenant $T'$, qui est l'image inverse d'un sous-groupe de $W'=N(T')/T'$ qui est \deleted{\unsure{étale, de rang|étale d'ordre} \inserted{fini,} $h$ sur $S$, et à fibres cycliques}\inserted{can. isom. à} $\mu_h$ (donc $h$ inv. sur $S$) et dont un générateur correspond à un élément de Coxeter.
% archive sheet number: 47

